\section{A practical method for producing theorems}\label{sec:method}

This section turns the pipeline of Figure~\ref{fig:pipeline} into a recipe that a student, or a program, can follow. Each step has an input, an output, and a check.

\subsection{Step 1: choose an invariant you can read on both sides}

Start from the arch registry. For presheaf topoi, each arch comes with a reading on the shape $C$ and an \emph{unfolding rule} that turns the property of the topos into a sentence about the structures:

\begin{center}\small
\begin{tabularx}{\linewidth}{>{\raggedright\arraybackslash}p{2.6cm}>{\raggedright\arraybackslash}p{3.2cm}>{\raggedright\arraybackslash}X>{\raggedright\arraybackslash}X}
\toprule
invariant & reading on $C$ & on right $M$-sets & on representations of a quiver\\
\midrule
Boolean & groupoid & $M$ is a group & no arrows\\
De Morgan & Ore condition & $aM\cap bM\neq\emptyset$ & $\le 1$ outgoing arrow per vertex\\
Morita class & $\Kar(C)$ & corners $eMe$, $1\in MeM$ & the quiver itself\\
essential points & objects of $\Kar(C)$ & idempotents mod $\mathcal D$ & vertices\\
subtopoi & topologies on $C$ & idempotent ideals ($M$ finite) & e.g.\ $2^n+1$ for the $n$-cycle\\
\bottomrule
\end{tabularx}
\end{center}

\emph{Check:} the invariant must be invariant under equivalence of categories. Properties stated with subobjects, limits and colimits are; properties that mention a chosen presentation (\q{the shape has three objects}) are not.

\subsection{Step 2: recognise a deck}

A structure is a presheaf topos when it has the following anatomy:
\begin{enumerate}
\item a set of data at each \q{place} (an object of $C$);
\item a way to transport data along \q{relations} between places (restriction maps), contravariant and compatible with composition;
\item no further equations or gluing conditions.
\end{enumerate}
If (3) fails because data must glue (functions on open sets, sections of a covering), the structure is a sheaf topos on a site. Bridge~\ref{br:sub} still applies there; the other four arches need versions for sites. Typical decks: actions of a monoid ($C = \B M$), representations of a quiver ($C$ = free category), simplicial objects ($C = \Delta$), families indexed by a preorder ($C = P^{\op}$), models of a Horn theory ($C$ = finitely presented models, opposite). \emph{Check:} write down $C$ explicitly and verify, on one or two examples, that a presheaf on $C$ is the structure you had in mind; the side of the action and the direction of arrows are the usual sources of error.

\subsection{Step 3: specialise the arch}

Compute the reading on $C$. Two heuristics cover most cases:
\begin{itemize}
\item \emph{Test on representables.} A property of all presheaves is usually decided on the representable presheaves $\y U = \Hom(-,U)$ and their sieves; this is how Bridges~\ref{br:bool} and~\ref{br:dm} are proved, and how counterexamples are built (Figure~\ref{fig:branch}).
\item \emph{Use a normal form for arrows.} In a free category arrows are paths, in $\B M$ elements, in a preorder there is at most one between two objects, so the condition becomes a statement about words, elements or inequalities.
\end{itemize}
\emph{Check:} enumerate small cases by machine when possible (small monoids, small quivers), and compare with the specialised condition.

\subsection{Step 4: unfold and rewrite}

Translate the property of the topos with the unfolding rule (Lemma~\ref{lem:neg} for negation), then translate the condition on $C$ into the usual vocabulary of the field (minimal right ideals, out-degree, rank). Keep both forms: the raw form is what Lean proves, the rewritten form is what a reader understands.

\subsection{Step 5: verify twice}

\emph{In Lean.} A generated statement is a two-line corollary of the arch once the condition on $C$ is proved. Here is the pattern, from the library:
\begin{leancode}
theorem isDeMorganTopos_singleObj_iff :
    IsDeMorganTopos ((SingleObj M)ᵒᵖ ⥤ Type) ↔ ∀ a b : M, ∃ c d : M, a * c = b * d := by
  rw [isDeMorganTopos_presheaf_iff_rightOreCondition, rightOreCondition_singleObj_iff]

theorem not_isDeMorganTopos_freeMonoid_two :
    ¬ IsDeMorganTopos ((SingleObj (FreeMonoid (Fin 2)))ᵒᵖ ⥤ Type) := by
  rw [isDeMorganTopos_singleObj_iff]
  -- remains: x and y have no common right multiple in the free monoid
\end{leancode}
All the work is in the lemma on $C$ (here \lean{rightOreCondition_singleObj_iff}), which is ordinary algebra or combinatorics. \emph{By hand.} Write a direct proof that does not use topoi; if it is much harder than the bridge, the bridge has bought something.

\subsection{Step 6: filter for novelty}

Write a one-sentence slogan for the Lean form of the rewritten statement and look it up in the TheoremGraph matching dataset \cite{theoremgraph-matching}; describe the statement in natural language to the semantic search of TheoremSearch \cite{theoremsearch}; look up numerical outputs in the OEIS \cite{oeis}. Then search the literature by hand, both for the rewritten statement in the vocabulary of its field and for the topos-theoretic one. Statements found in print go to the calibration log (\S\ref{sec:calibration}). A miss shows nothing: \q{not found} is not \q{new}. Appendix~\ref{app:novelty} gives each pass, its parameters and its results.

\subsection{Step 7: decide whether it is interesting}

A generated statement is worth keeping when (a) it connects two vocabularies that are not usually connected, (b) its rewritten form is short, (c) its direct proof is not a one-liner \emph{or} the question it answers had not been asked. We drop a candidate that fails all three. We did not record how many were dropped (\S\ref{sec:calibration}), and ranking the candidates is one of the limits of the method (\S\ref{sec:limits}).
